% Prologue of The Hydrogen Atoms of Deep Learning.
\specialchapter{Prologue}{The hydrogen atom, and what a solvable model is for}{4}\label{ch:prologue}
\begin{refsection}[chapters/00-prologue.bib]

Every physics student solves the same three problems before solving anything else: the harmonic oscillator, the hydrogen atom, and the Ising model. None of them describes a real system exactly, and every student learns, along with the solutions, the specific ways in which each one is wrong. This book proposes the same education for deep learning. It collects the models of deep learning that can be solved completely, solves each one in a single notation, and says, for each, exactly where it fails.

\addsec{Three models every physicist solves}

\paragraph{The hydrogen atom.} In 1885 Johann Balmer, a teacher at a girls' school in Basel, found that the wavelengths of the four visible lines of hydrogen fit the formula $\lambda = B\,n^2/(n^2-4)$ for $n=3,4,5,6$, with a single constant $B\approx364.56$ nanometers \parencite{Balmer1885}. His formula was a regularity without a reason. Johannes Rydberg generalized it to other elements and wrote it in terms of wave numbers, $1/\lambda = R(1/4-1/n^2)$ \parencite{Rydberg1890}. In 1913 Niels Bohr derived Rydberg's constant from a model in which the electron's angular momentum is quantized \parencite{Bohr1913}, and in 1926 Erwin Schr\"odinger derived the whole spectrum as the eigenvalues of a differential operator, in the first paper on his wave equation \parencite{Schrodinger1926}. The lines on the title page of this book stand at the wavelengths that Rydberg's formula predicts.\sidenote{The emblem draws the Balmer series, $n=3,4,\dots$, with the four visible lines labelled; the dotted line marks the series limit near $364.6$ nanometers, where the lines accumulate.}

The solution is exact, and it is wrong. Arnold Sommerfeld had already found in 1916 that relativistic corrections split the lines into closely spaced multiplets, the \term{fine structure}\index{fine structure} \parencite{Sommerfeld1916}; Paul Dirac's relativistic equation of 1928 accounted for it \parencite{Dirac1928}. Dirac's theory in turn predicted that two levels of the $n=2$ shell coincide. In 1947 Willis Lamb and Robert Retherford measured a small splitting between them \parencite{LambRetherford1947}, and within weeks Hans Bethe explained most of it by the interaction of the electron with the quantized electromagnetic field \parencite{Bethe1947}. The Lamb shift became one of the starting points of quantum electrodynamics. Each correction to the hydrogen atom names a piece of physics that the previous model had left out: relativity, then the field.

\paragraph{The Ising model.} Wilhelm Lenz proposed a lattice of spins that can point up or down and prefer to align with their neighbors; his student Ernst Ising solved it in one dimension in 1925, found no phase transition, and conjectured that there was none in higher dimensions either \parencite{Ising1925}. The conjecture was false. Rudolf Peierls argued in 1936 that the two-dimensional model does order at low temperature \parencite{Peierls1936}, and Lars Onsager computed its free energy exactly in 1944 \parencite{Onsager1944}. The cheaper solution, mean-field theory, is exact when every spin interacts with every other, and it predicts that the magnetization vanishes at the critical temperature like $(T_c-T)^{1/2}$. In two dimensions the exponent is $1/8$, as C.\,N.\ Yang derived from Onsager's solution \parencite{Yang1952}. Mean-field theory gets the phase diagram qualitatively right and the critical behavior wrong, and the reason it is wrong (fluctuations matter in low dimension) became the subject of the renormalization group.

\paragraph{The harmonic oscillator.} Expand a potential to second order around a minimum and every system becomes a set of independent oscillators, its normal modes. The approximation is the most useful in physics and it is wrong in instructive ways. A perfectly harmonic crystal does not expand when heated, and heat flows through it without resistance, because its normal modes never exchange energy; both thermal expansion and finite thermal conductivity come from the anharmonic terms that the expansion dropped \parencite{AshcroftMermin1976}.

\medskip

What these three have in common is not that they are true. Each is \emph{solvable}, because a symmetry or a limit makes it so (rotational invariance for hydrogen, a transfer matrix for Ising, quadratic energy for the oscillator). Each reproduces something real. Each makes quantitative predictions. And each fails in ways that are known, specific, and informative, so that learning the failure is learning the next piece of physics. George Box's remark that all models are wrong is often quoted as an excuse for a model \parencite{Box1979}. The physics curriculum treats it as an instruction: say how.

\addsec{The models of deep learning}

Deep learning has no standard model, and it may never have one: a trained network is a large, specific object shaped by its data, its architecture and its training procedure, and the phenomena it displays are many. But the theory of deep learning has, over the last decade, accumulated a surprising number of small models that can be solved exactly and that reproduce one phenomenon each. Least squares with Gaussian features shows a test error that falls, rises and falls again. A linear network with two layers learns in stages. A network whose weights are products of two numbers prefers sparse solutions when started small. Gradient descent on a single neuron in high dimension spends most of its time searching and then learns quickly. A small network trained on modular addition computes with Fourier modes. A linear attention layer trained on regression prompts performs one step of gradient descent on the prompt. A product of two scalars trained with a large step size jumps out of a sharp valley into a flat one.

These models are scattered across dozens of papers, written in different notations, by communities that rarely cite each other: statistical physics, high-dimensional statistics, optimization, theoretical computer science, and mechanistic interpretability. The places where the models fail are written down even less often, and are usually in other papers. A student who wants the deep learning analogue of the hydrogen atom has to assemble it from the papers, and then discover, unaided, what it gets wrong.

This book assembles sixteen of them. Each chapter takes one phenomenon observed in trained networks, gives the smallest model that reproduces it, proves the model's main theorem completely, and then says precisely what the model gets wrong, with references to the results that show it.\sidenote{The book is not a survey. Many important results of deep learning theory are absent because they are not solvable models of a phenomenon: generalization bounds, approximation theory, most of optimization theory. The monograph of \textcite{Zhao2026} covers the field's breadth, one topic per chapter.}

\addsec{What makes a model a hydrogen atom}

We hold every model in the book to four tests.

\begin{enumerate}
\item \emph{It reproduces a phenomenon.} Something observed in real networks, robustly enough that we can state it as an observation with evidence, appears in the model.
\item \emph{It is solvable.} Its main claim has a complete proof that a graduate student can check, short enough to teach in a lecture or two.
\item \emph{It predicts a number.} An exponent, a threshold, a time scale, a phase boundary, or a constant, which can be compared with an experiment.
\item \emph{It fails instructively.} The ways in which it is wrong are known and specific, and each failure points to a missing ingredient and to the next model.
\end{enumerate}

The tests exclude a great deal. A model that reproduces nothing observed is mathematics, and may be good mathematics, but it is not a model of deep learning. A heuristic that cannot be made rigorous may be right, and when it is the best available we include it, labelled as such; but the core of every chapter is a theorem. A model that predicts nothing quantitative cannot be tested, and a model whose failures are unknown has not been tested. The fourth test is the one most often skipped, and it is the one this book takes most seriously.

\addsec{A first hydrogen atom: gradient descent on a quadratic}

The simplest solvable model of training is the harmonic oscillator of optimization: gradient descent on a quadratic loss. It is worth solving carefully, because almost every chapter of this book is about one of its failures.

Let $H$ be a symmetric positive semidefinite $p\times p$ matrix with eigenvalues $\lambda_1\ge\dots\ge\lambda_p\ge0$ and orthonormal eigenvectors $e_1,\dots,e_p$, and let
\[
  \Loss(\theta)=\tfrac12\,(\theta-\theta_\star)^\top H(\theta-\theta_\star).
\]
Gradient descent with step size $\eta>0$ is the iteration $\theta_{t+1}=\theta_t-\eta\nabla\Loss(\theta_t)=\theta_t-\eta H(\theta_t-\theta_\star)$. Least squares is the main example: if $X\in\R^{n\times p}$ and $y\in\R^n$, then $\frac1{2n}\norm{X\theta-y}^2$ equals a quadratic of this form plus a constant, with $H=X^\top X/n$ and $\theta_\star$ a least-squares solution.

\begin{theorem*}[Normal modes of gradient descent]
For every initialization $\theta_0$ and every $t\ge0$, the error $\theta_t-\theta_\star$ has coordinates
\[
  \ip{\theta_t-\theta_\star,e_k}=(1-\eta\lambda_k)^t\,\ip{\theta_0-\theta_\star,e_k}\qquad(k=1,\dots,p),
\]
and the loss is $\Loss(\theta_t)=\frac12\sum_k\lambda_k(1-\eta\lambda_k)^{2t}\ip{\theta_0-\theta_\star,e_k}^2$. Consequently:
\begin{enumerate}
\item The component along $e_k$ converges to zero for every initialization if and only if $0<\eta\lambda_k<2$; it stays constant if $\lambda_k=0$, and it grows without bound, with alternating sign, if $\eta\lambda_k>2$ and it is nonzero initially.
\item If $H$ is positive definite, gradient descent converges for every initialization if and only if $\eta<2/\lambda_1$, and the step size $\eta=2/(\lambda_1+\lambda_p)$ gives $\norm{\theta_t-\theta_\star}\le\big(\frac{\kappa-1}{\kappa+1}\big)^t\norm{\theta_0-\theta_\star}$ with $\kappa=\lambda_1/\lambda_p$.
\end{enumerate}
\end{theorem*}

\begin{proof}
The error $\delta_t=\theta_t-\theta_\star$ satisfies $\delta_{t+1}=(I-\eta H)\delta_t$, so $\delta_t=(I-\eta H)^t\delta_0$. The matrix $I-\eta H$ has the eigenvectors $e_k$ with eigenvalues $1-\eta\lambda_k$, which gives the first formula, and the loss is $\frac12\sum_k\lambda_k\ip{\delta_t,e_k}^2$. A geometric sequence $(1-\eta\lambda_k)^tc$ with $c\ne0$ tends to zero exactly when $|1-\eta\lambda_k|<1$, that is, when $0<\eta\lambda_k<2$; it is constant when $\lambda_k=0$, and when $\eta\lambda_k>2$ its ratio is below $-1$, which is part (i). For part (ii), all $\lambda_k$ are positive, so (i) applies to every component, and the largest eigenvalue gives the binding constraint. With $\eta=2/(\lambda_1+\lambda_p)$ every factor satisfies $|1-\eta\lambda_k|\le(\lambda_1-\lambda_p)/(\lambda_1+\lambda_p)=(\kappa-1)/(\kappa+1)$, because $1-\eta\lambda_k$ decreases in $\lambda_k$ from $(\lambda_1-\lambda_p)/(\lambda_1+\lambda_p)$ at $\lambda_k=\lambda_p$ to its negative at $\lambda_k=\lambda_1$; orthogonality of the $e_k$ gives the bound on $\norm{\delta_t}$.
\end{proof}

Three predictions follow, and each will return.

\emph{A stability threshold.} The step size must stay below $2/\lambda_1$, twice the inverse of the largest curvature. Above it, the stiffest mode oscillates with growing amplitude (\cref{fig:prologue:gd}).

\emph{Learning by modes.} Different directions are learned at different speeds. In the limit of small steps, gradient flow $\dot\theta=-H(\theta-\theta_\star)$ gives $\ip{\theta(t)-\theta_\star,e_k}=e^{-\lambda_kt}\ip{\theta(0)-\theta_\star,e_k}$, so the mode $e_k$ is learned on the time scale $1/\lambda_k$: the large-eigenvalue directions first and the small ones last. In kernel regression the large-eigenvalue directions are the smooth functions, and this simple fact becomes the \emph{spectral bias} of \cref{ch:lazy}.

\emph{An implicit choice.} When $H$ is singular, the minimizers form an affine subspace, and gradient descent picks one of them without being told which. The components of the error along the null space of $H$ never change, and when $\eta$ is below $2/\lambda_1$ the other components tend to zero, so the limit is the minimizer closest to the initialization. For least squares with $n<p$, started at $\theta_0=0$, every iterate lies in the row space of $X$, because every gradient $X^\top(X\theta-y)/n$ does; the limit is therefore the minimum-norm interpolator $X^+y$. This interpolator is the protagonist of \cref{ch:doubledescent}, and \cref{ch:diagonal} shows how changing the parameterization changes the choice.

\begin{figure}[t]
\centering
\begin{fullwidth}
\centering
\begin{tikzpicture}
\begin{groupplot}[group style={group size=2 by 1,horizontal sep=16mm},hplot,width=80mm,height=62mm]
\nextgroupplot[xlabel={$\theta_1$},ylabel={$\theta_2$},xmin=-2.9,xmax=0.6,ymin=-1.15,ymax=1.15]
\pgfplotsinvokeforeach{0.3,1.2,2.7,4.9}{\addplot[seamist!80!sea,line width=0.4pt,domain=0:360,samples=90] ({sqrt(2*#1)*cos(x)},{sqrt(2*#1/10)*sin(x)});}
\addplot[sea,line width=0.9pt,mark=*,mark size=0.9pt] table[x=x,y=y] {figures/data/00-prologue-gd-a.dat};
\addplot[sandline!55!black,line width=0.8pt,mark=*,mark size=0.9pt] table[x=x,y=y,restrict expr to domain={\thisrow{t}}{0:14}] {figures/data/00-prologue-gd-b.dat};
\addplot[seadeep,line width=0.7pt,densely dotted,mark=*,mark size=0.9pt] table[x=x,y=y,restrict expr to domain={\thisrow{t}}{0:5}] {figures/data/00-prologue-gd-c.dat};
\nextgroupplot[xlabel={step $t$},ylabel={loss $\Loss(\theta_t)$},ymode=log,xmin=0,xmax=40,ymin=1e-7,ymax=1e3,legend pos=south west]
\addplot[theory] table[x=t,y=exact] {figures/data/00-prologue-gd-a.dat};
\addplot[theory3] table[x=t,y=exact] {figures/data/00-prologue-gd-b.dat};
\addplot[theory2] table[x=t,y=exact] {figures/data/00-prologue-gd-c.dat};
\addplot[simulation,each nth point=3] table[x=t,y=loss] {figures/data/00-prologue-gd-a.dat};
\addplot[simulation,each nth point=3] table[x=t,y=loss] {figures/data/00-prologue-gd-b.dat};
\addplot[simulation,each nth point=3] table[x=t,y=loss] {figures/data/00-prologue-gd-c.dat};
\legend{$\eta=0.05$,$\eta=0.17$,$\eta=0.205$}
\end{groupplot}
\end{tikzpicture}
\end{fullwidth}
\caption{Gradient descent on $\Loss(\theta)=\frac12(\theta_1^2+10\,\theta_2^2)$, for which $2/\lambda_1=0.2$. Left: iterates from $\theta_0=(-2.6,0.55)$ over level sets of the loss; with $\eta=0.05$ the path is smooth, with $\eta=0.17$ the stiff coordinate flips sign at every step but shrinks, and with $\eta=0.205$ it flips and grows (first five steps shown). Right: the loss, computed by running gradient descent (points) and from the normal-mode formula (curves); they agree to machine precision, which is what ``solvable'' means.}
\label{fig:prologue:gd}
\end{figure}

\addsec{What the quadratic gets wrong}

The quadratic is exact for linear regression. As a model of training a network it fails in at least five ways, and the list of its failures is close to a table of contents for this book.

\paragraph{The curvature is not constant.} The Hessian of a network's loss changes during training. The theorem forbids step sizes above $2/\lambda_1$, but networks are routinely trained with step sizes for which the curvature at initialization, or later in training, exceeds $2/\eta$, and the training does not diverge: the curvature adapts to the step size. A model with only two parameters already shows how (\cref{ch:stability}).

\paragraph{The loss is not convex.} A network's loss has saddle points, where the gradient vanishes but the loss is not minimal, and training slows down near them. A quadratic has no saddles, so it cannot produce the plateaus of training curves. Escaping from saddles is the subject of the chapters on deep linear networks (\cref{ch:deeplinear}), diagonal networks (\cref{ch:diagonal}), single-index models (\cref{ch:singleindex}), committee machines (\cref{ch:committee}), sparse parities (\cref{ch:parity}) and induction heads (\cref{ch:induction}).

\paragraph{The minimizer is not unique.} A network that can interpolate its training data has a manifold of minimizers, and which one training finds is decided by the parameterization, the initialization and the optimizer. In the quadratic this choice is the trivial projection described above. In networks it produces sparsity (\cref{ch:diagonal}), maximum margins (\cref{ch:margin}), Fourier features (\cref{ch:grokking}) and the simplex geometry of neural collapse (\cref{ch:collapse}).

\paragraph{The data have no structure.} In least squares the matrix $H$ is the covariance of the inputs, and its spectrum is all the quadratic knows about the data. That is enough for some purposes: the singular values of a random design produce double descent (\cref{ch:doubledescent}), and power-law spectra produce power-law learning curves (\cref{ch:scaling}). It is not enough for others: a target that depends on one hidden direction of a high-dimensional input (\cref{ch:singleindex}), or on the parity of a few bits (\cref{ch:parity}), is invisible to second-order statistics.

\paragraph{There are no features.} A quadratic in the parameters is a linear model, and a linear model has fixed features. The central fact about trained networks is that their features change. Wide networks can be made to behave like linear models (\cref{ch:lazy}), and the price of doing so is a provable failure to learn targets that networks with moving features learn easily. The chapters on what networks represent (\cref{ch:grokking,ch:superposition,ch:collapse}) and on transformers (\cref{ch:icl,ch:induction}) are about features that training builds.

\addsec{How each chapter is built}

Every chapter after this one has the same anatomy.

\emph{The phenomenon} states what is observed in trained networks, by whom and with what evidence, in an \emph{observation} environment that separates empirical regularities from theorems, and ends with the question the model will answer.

\emph{The model} defines the smallest solvable model of the phenomenon, says what was thrown away and why, and gives a \emph{dictionary} that translates the model's quantities into the quantities of a trained network.

\emph{The solution} proves the main theorem completely. Further results come with complete proofs where they fit, with proof sketches that say what is omitted, or with heuristic derivations. The three are marked differently: a proof ends with $\blacksquare$, a sketch is labelled as one, and a heuristic derivation, the kind of argument that physicists trust and mathematicians do not yet accept, ends with $\Diamond$. The difference matters. A heuristic derivation can be right and useful, and several of the most informative predictions in this book are derived rather than proved, but a reader should always know which is which.

\emph{Reading the solution} turns the theorem into predictions and compares them with simulations. Every chapter has at least one figure in which a curve is the model's exact prediction and the points are a simulation, produced by a script that is distributed with the book.

\emph{What the model gets wrong} has three parts: what was thrown away, where the model breaks (specific theorems or experiments that contradict it), and the fine structure (the corrections and successor models that repair it, and what is still open).

A \emph{model card} summarizes the chapter in six lines: phenomenon, model, solution, prediction, failure and successors. The exercises come in four kinds: \emph{consolidation} exercises check and extend the chapter's results; \emph{variation} exercises change one assumption of the model and ask for the new solution, which is the skill this book most wants to teach; \emph{simulation} exercises ask for a numerical experiment and say what it should show; \emph{research} exercises point to open problems. The notes at the end of each chapter close with a paragraph that translates the notation of the main source papers into the notation of this book.

\addsec{A map of the book}

\Cref{tab:prologue:map} lists the models. The six parts move from models that are linear in their parameters (Part I), through the two infinite-width limits (Part II), to models in which features are learned (Part III) and in which the learned structure can be read off the weights (Part IV), then to transformers (Part V) and to large step sizes and large models (Part VI). The epilogue looks across the chapters at what the models share and at what they get wrong together, and lists phenomena that still lack a hydrogen atom.

\begin{table}[p]
\begin{fullwidth}
\centering\small
\renewcommand{\arraystretch}{1.18}
\begin{tabular}{@{}r>{\raggedright\arraybackslash}p{28mm}>{\raggedright\arraybackslash}p{33mm}>{\raggedright\arraybackslash}p{41mm}>{\raggedright\arraybackslash}p{38mm}@{}}
\toprule
& \textsc{phenomenon} & \textsc{model} & \textsc{solution} & \textsc{gets wrong}\\
\midrule
1 & Double descent & Least squares with Gaussian features & Exact risk at every $n$ and $p$; ridge in the proportional limit & Features are fixed and isotropic\\
2 & Learning in stages & Deep linear networks & Exact decoupled mode dynamics, logistic for two layers & Only linear maps; aligned initialization\\
3 & Implicit bias of small initialization & Diagonal linear networks & A mirror flow interpolating between $\ell_2$ and $\ell_1$ & A linear predictor in disguise\\
4 & Training past zero error & Linear classifiers, exponential loss & The margin tends to the maximum margin & Logarithmic rates; networks reach only KKT points\\
\midrule
5 & Depth needs critical initialization & Wide random networks & An order--chaos phase diagram & Finite width matters when depth is comparable\\
6 & Wide networks as kernels & Linearization at initialization & Lazy training; a lower bound for every kernel & No feature learning\\
\midrule
7 & A search phase before learning & Single-index models & Online SGD needs about $d^{k-1}$ samples & Reusing data beats the exponent\\
8 & Plateaus and specialization & The soft committee machine & Exact ODEs; a symmetric saddle & Gaussian inputs and tiny steps\\
9 & Abrupt learning, hidden progress & Sparse parities & A statistical-query bound; one gradient step finds the support & Natural tasks have low-degree signal\\
\midrule
10 & Grokking & Modular addition & A Fourier network; margin-maximizing neurons are Fourier modes & The mechanism varies between runs\\
11 & Superposition & A ReLU toy autoencoder & Linear models never superpose; sparse features force it & Features are given and independent\\
12 & Neural collapse & Unconstrained features & Minimizers are simplex frames & Features are not free\\
\midrule
13 & In-context learning & Linear attention on regression & The optimum is one preconditioned gradient step & No softmax; Gaussian tasks\\
14 & Induction heads & Attention on random Markov chains & Bayes predictor; a two-layer circuit; a saddle & One circuit among many\\
\midrule
15 & The edge of stability & The product $uv$ & An exact law for the imbalance forces the catapult & No progressive sharpening\\
16 & Scaling laws & Power-law spectra and Zipf's law & Exponents from tails; compute-optimal allocation & The exponents are inputs\\
\bottomrule
\end{tabular}
\caption{The sixteen models, each with the phenomenon it reproduces, the main result proved in its chapter, and its most important failure.}
\label{tab:prologue:map}
\end{fullwidth}
\end{table}

\addsec{How to read this book}

The prerequisites are undergraduate analysis, linear algebra and probability, and familiarity with gradient descent and with neural networks at the level of an introductory course. The chapters are written to be read independently, and each builds what it uses; tools that several chapters share, such as Hermite polynomials, the Marchenko--Pastur law, martingale inequalities and the method of ordinary differential equations for stochastic recursions, are collected in \cref{ch:toolkit}. The notation is uniform across the book and is summarized before this prologue. A one-semester course can take Parts I through III in order; a course on interpretability can start with Part IV and continue with Part V; a reader interested in optimization can read this prologue and \cref{ch:deeplinear,ch:diagonal,ch:margin,ch:stability}.

Three cautions. First, a solvable model is not an explanation of deep learning. It is an explanation of one phenomenon in one idealized setting, and the claim that the same mechanism operates in a trained network is an empirical claim that each chapter examines separately. Second, the models in this book were chosen because they are solvable, and solvable models are biased toward phenomena that have simple mechanisms; the epilogue discusses what this bias hides. Third, the field moves quickly. The models here are the ones that have earned a place by 2026, and several chapters end with successor models that may soon deserve chapters of their own.

\notesandrefs

The history of the hydrogen spectrum, from Balmer's formula to quantum electrodynamics, is told in most textbooks of atomic physics; Balmer's formula, Bohr's model and Schr\"odinger's solution are in the original papers cited above, and Lamb and Retherford's measurement and Bethe's calculation appeared in the same volume of the \emph{Physical Review} in 1947. Onsager's solution of the two-dimensional Ising model \parencite{Onsager1944} and Yang's computation of its spontaneous magnetization \parencite{Yang1952} are the classic exact results; the failure of mean-field exponents in low dimension is the starting point of every account of critical phenomena. The role of anharmonicity in thermal expansion and thermal conductivity is in \textcite{AshcroftMermin1976}.

The convergence theory of gradient descent on quadratics, including the optimal constant step size, is standard; see \textcite{Polyak1987} and \textcite{Nesterov2018}. The observation that gradient descent on underdetermined least squares from zero converges to the minimum-norm solution is folklore in numerical analysis and became central to deep learning theory through the study of implicit regularization.

Four books are close to this one in spirit. \Textcite{EngelVandenBroeck2001} present the statistical mechanics of learning, including several models that reappear here in modern form. \Textcite{RobertsYaidaHanin2022} develop a single method, an effective field theory of wide networks organized as an expansion in the inverse width, and carry it much further than a single chapter could. \Textcite{Huang2021} covers the classical statistical-physics models of neural networks, such as spin glasses, the Hopfield model and the perceptron. \Textcite{Zhao2026}, a monograph posted on arXiv, gives a broad, proof-oriented account of deep learning theory in which each topic has about one chapter. The position paper of \textcite{SimonEtAl2026} argues that a scientific theory of deep learning is emerging and names solvable idealized settings as one of its five strands; this book is a textbook for that strand.

\printchapterbib
\end{refsection}
